\documentclass[conference]{IEEEtran}
\usepackage[utf8]{inputenc}
\usepackage[T1]{fontenc}
\usepackage[turkish]{babel}
\usepackage{graphicx}
\usepackage{float}
\usepackage{caption}
\usepackage{listings}
\usepackage{algorithm}
\usepackage{algpseudocode}
\usepackage{xurl}
\usepackage[hidelinks]{hyperref}

\newcommand{\LongState}[1]{%
	\State \parbox[t]{\dimexpr\linewidth-\algorithmicindent\relax}{#1}%
}

% Görselleri sütun genişliğinin %80'i kadar ölçekle ve en-boy oranını koru
\setkeys{Gin}{width=0.8\columnwidth,keepaspectratio}

\makeatletter
\renewenvironment{abstract}{%
	\par\noindent
	{\normalfont\normalsize ÖZET}\par\medskip
	\normalfont\normalsize
}{%
	\par
}
\makeatother

\title{%
	\normalsize\bfseries
	KOCAELİ ÜNİVERSİTESİ BİLGİSAYAR MÜHENDİSLİĞİ\\[1em]
	YAZILIM LABORATUVARI - II\\
	PROJE - II\\[1em]
	Word Crush Mobil Oyunu
}

\author{
	\IEEEauthorblockN{Saffet Hakan KOÇAK}
	\IEEEauthorblockA{%
		Bilgisayar Mühendisliği\\
		Kocaeli Üniversitesi\\
		230202058%
	}
	\and
	\IEEEauthorblockN{Yusuf BÜLBÜL}
	\IEEEauthorblockA{%
		Bilgisayar Mühendisliği\\
		Kocaeli Üniversitesi\\
		230202050%
	}
}

\begin{document}
\maketitle

\begin{abstract}
Bu projede, Android tabanlı Word Crush mobil oyunu tasarlanmış ve
geliştirilmiştir. Oyunun amacı, oyuncunun harflerden oluşan bir grid
üzerinde anlamlı kelimeler üretmesini sağlamak; oluşan kelimeleri tahtadan
kaldırıp yeni harf akışlarıyla oynanışı sürekli ve dengeli tutmaktır. Proje,
kelime doğrulama ve puanlama kurallarının tutarlı biçimde uygulanması ile
mobil cihazlarda akıcı bir oyun deneyimi sunmayı hedeflemektedir.

Uygulama mimarisi MVVM yapısında kurgulanmış; arayüz Jetpack Compose ile
geliştirilmiş, veri katmanında Room ve DataStore kullanılmıştır. Oyun
durumu, skor tablosu, kullanıcı adı ve ilerleme bilgileri kalıcı olarak
saklanırken; iş kuralları use-case odaklı bir katmanda ayrılmış ve
görünüm-modeli ile arayüz arasında tek yönlü veri akışı sağlanmıştır.

Oyun mekanikleri; kelime oluşturma, geçersiz hamle engelleme, kelime
silme-sonrası harf düşürme, yeni harf üretimi, combo hesaplama ve
puanlama kuralları üzerinden tanımlanmıştır. Joker sistemi, market ve
özel güç mekanikleri ile stratejik kararların önemi artırılmış; skor
tablosu ve kullanıcı adı saklama özellikleri ile rekabet duygusu
desteklenmiştir.

Geliştirilen yapıda oyun akışı, arayüz tepkiselliği ve veri kalıcılığı
uçtan uca test edilmiş; temel senaryolarda stabil bir oynanış elde
edilmiştir. Sonuç olarak, kelime tabanlı bulmaca dinamiklerini mobil
ortama uygun şekilde birleştiren, modüler ve genişletilebilir bir oyun
altyapısı ortaya konulmuştur.
\end{abstract}

\section{GİRİŞ}

Mobil oyunlar, geniş kullanıcı kitlesine erişimi, kısa ve tekrarlı
etkileşimlere uygun yapısı ve cihazların her an kullanılabilir olması
nedeniyle günümüzde etkin bir dijital deneyim türü haline gelmiştir.
Kelime tabanlı oyunlar ise eğitsel katkı ile eğlenceyi birleştirerek
kullanıcıya dilsel farkındalık, kelime haznesini geliştirme ve dikkat
yönetimi gibi bilişsel faydalar sunar. Bu bağlamda kelime oluşturma ve
bulmaca mekaniklerini mobil platformlara uygun şekilde bir araya getirmek,
hem kullanışlı hem de sürdürülebilir bir oyun deneyimi açısından önemlidir.

Word Crush projesi, mobil ortamda akıcı bir kelime oyunu deneyimi sunmak
üzere geliştirilmiştir. Projede hedeflenen, oyuncunun harflerden oluşan
bir grid üzerinden anlamlı kelimeler üretmesini sağlayan; puanlama, combo,
joker ve özel güç mekanikleri ile stratejik derinlik kazandıran bir oyun
altyapısıdır. Oyuncuya sunduğu mobil deneyim, hem tek elle hızlı
etkileşim hem de performanslı arayüz güncellemeleriyle desteklenmiştir.

Oyun yapısı; grid tabanlı tahtanın yönetimi, kelime doğrulama kuralları,
kelime silinmesi sonrasında harflerin düşürülmesi ve yeni harf üretimi
akışları üzerine kuruludur. Skor ve combo mekanikleri puanlama dengesi
için tanımlanmış; joker sistemi ve market yapısı oyuncunun oyun içindeki
kaynaklarını yönetmesine olanak vermiştir. Skor tablosu, kullanıcı adı
saklama ve oyun durumu kaydı gibi bileşenler veri kalıcılığını sağlayarak
uzun süreli kullanımı destekler.

Raporun devamında, uygulamanın mimarisi ve katmanları, oyun mekaniklerinin
algoritmik işleyişi, veri saklama yaklaşımları ve arayüz tasarımı
ayrıntılı olarak ele alınacak; ardından deneysel sonuçlar, elde edilen
bulgular ve gelecek iyileştirme alanları sunulacaktır.

\section{Yöntem}

Bu bölümde, Word Crush oyununun tasarım ve uygulama yaklaşımı, mimari
katmanlar, oyun akışı ve temel mekanikler üzerinden açıklanmaktadır.
Sistem; arayüz, alan mantığı ve veri saklama katmanlarını ayıran bir
yapı üzerinde kurulmuş; oyun kurgusu ise grid tabanlı etkileşim ve
kelime doğrulama kuralları etrafında şekillendirilmiştir.

\subsection{Sistem Mimarisi}

Uygulama, MVVM mimarisi temelinde katmanlara ayrılarak
geliştirilmiştir. UI katmanı Jetpack Compose üzerinden yalnızca
görünümü ve etkileşimleri yönetir; ViewModel katmanı, oyun
durumunu tek yönlü veri akışı ile taşır ve UI ile iş kuralları
arasındaki bağlantıyı kurar. Veri katmanı, Room veritabanı
üzerinden skor tablosu ve oyun kayıtlarını saklarken; DataStore
ile kullanıcı adı ve ayar gibi küçük durum bilgilerini kalıcı
hale getirir. Oyun motoru katmanı ise grid durumu, kelime
doğrulama kuralları, puanlama, combo ve özel güç mekaniklerini
yöneterek hamlelerin tutarlı şekilde işlenmesini sağlar.
Ekranlar arası akış, navigation bileşenleri ile kontrol edilmiş
ve oyun, market ve skor tablosu gibi görünümler arasında
durum geçişleri standartlaştırılmıştır.

\begin{figure}[htbp]
	\centering
	\includegraphics{sistem_mimari.png}
	\caption{Sistem Mimarisi Diyagramı}
	\label{fig:mimari}
\end{figure}

\subsection{Kullanılan Teknolojiler}

Uygulama Kotlin ile geliştirilmiş; arayüz katmanı Jetpack Compose ve
navigation bileşenleri ile uygulanmıştır. Kalıcı veri saklama için
Room veritabanı, küçük ayar ve kimlik bilgileri için DataStore
kullanılmıştır. ViewModel, oyun durumunu ekrana yansıtmak ve
durum geçişlerini yönetmek için merkezi rol üstlenmiştir.

\subsection{Oyun Akışı}

Oyun akışı, seviye başlatma, oyuncu hamlesi alma, kelime doğrulama,
puan hesaplama, tahta güncelleme ve oyun bitiş koşullarını
değerlendirme adımlarından oluşur. Bu akışta her hamle, oyun
durumunu günceller ve UI tarafında anlık geribildirim
üretilir; böylece akıcı ve izlenebilir bir oyun döngüsü sağlanır.

\begin{figure}[htbp]
	\centering
	\includegraphics{ana_ekran.png}
	\caption{Uygulama Ana Ekranı}
	\label{fig:ana_ekran}
\end{figure}

\subsection{Grid Yapısı ve Harf Üretimi}

Oyun tahtası, iki boyutlu kare bir grid üzerinde tanımlanmıştır
ve farklı zorluk seviyeleri için 6x6, 8x8 ve 10x10 seçenekleri
sunmaktadır. Her hücre, harf değeri, seçilebilirlik durumu ve
geçici işaretleme bilgilerini içeren bir veri modeliyle temsil
edilir; böylece oyun motoru, hamle kontrolü ve güncelleme
adımlarını hücre bazında yürütebilir.

Harf üretimi, Türkçe dilindeki harf frekansları dikkate alınarak
ağırlıklı bir dağılımla gerçekleştirilir. Sık kullanılan harfler
(örn. A, E, I, K, L, M, N, R), orta sıklıktaki harfler ve düşük
sıklıklı harfler (örn. J, F, G, H gibi) farklı olasılık katsayıları
ile üretilir. Bu yaklaşım, anlamlı kelime oluşumunu teşvik ederken
oyun zorluğunun kontrol edilebilir kalmasını sağlar.

Grid oluşturulduktan sonra, tahtada en az bir geçerli kelimenin
bulunabilirliği algoritmik olarak kontrol edilir. Kelime
üretilemeyen durumlarda tahta yeniden üretilir veya harfler
yeniden dağıtılır; böylece kullanıcıya çözülmez bir grid
gösterilmesi engellenir. Bu kontrol, oynanabilirlik garantisi
sağlayarak deneyimin kesintisiz ve adil ilerlemesine katkıda
bulunur.

\begin{figure}[htbp]
	\centering
	\includegraphics{yeni_oyun_ekrani.png}
	\caption{Yeni Oyun Ekranı}
	\label{fig:yeni_oyun}
\end{figure}

\subsection{Harf Seçme ve Kelime Oluşturma}

Oyuncu, grid üzerinde komşu harfleri ardışık olarak seçerek
kelime oluşturur. Komşuluk kuralı 8 yönlü olarak uygulanır;
yani yatay, dikey ve çapraz komşular seçim içine dahil
edilebilir. Aynı hücre, tek bir seçim dizisi içinde yeniden
kullanılamaz; bu sınır, kelime oluşturmayı tekil ve izlenebilir
bir rota üzerinden gerçekleştirir.

Seçim sırasında minimum 3 harf koşulu aranır ve bu koşul
sağlanmadan kelime finalize edilmez. Parmağın kaldırılması
veya seçim akışının tamamlanması ile kelime oluşturma
tamamlanır; oluşan dizi doğrulama adımına aktarılır ve
geçici seçim durumu sıfırlanır.

\begin{figure}[htbp]
	\centering
	\includegraphics{oyun_ekrani.png}
	\caption{Oyun Ekranı}
	\label{fig:oyun_ekrani}
\end{figure}

\subsection{Kelime Doğrulama Mekanizması}

Oluşan kelimeler, uygulama içindeki sözlük veri yapısı üzerinden
doğrulanır ve geçerlilik durumu belirlenir. Geçerli kelimelerde
puanlama ve tahta güncelleme akışı çalıştırılır; seçilen hücreler
temizlenir ve harf düşürme mekaniği devreye girer. Geçersiz
kelimelerde ise hamle etkisiz sayılır, seçim iptal edilir ve
oyuncuya geri bildirim sağlanır.

Hamle sayısı, hem geçerli hem de geçersiz kelimelerde azaltılır.
Bu tercih, stratejik seçim baskısını korur ve oyuncunun her
denemede dikkatli davranmasını sağlar. Doğrulama mekanizması,
oyunun dilsel tutarlılığını korurken hamle ekonomisini dengede
tutar.

\subsection{Puan Hesaplama}

Puanlama, harf bazlı bir sistem üzerinden hesaplanır ve her
harf için tanımlanan puan değerleri kelime puanını oluşturur.
Kelime puanı, seçilen harflerin tekil puanlarının toplanması
ile elde edilir; böylece kelimenin uzunluğu ve harf dağılımı
doğrudan skora yansır. Bu temel skora, combo ve özel güç
mekaniklerinden gelen çarpanlar eklenerek hamlelerin etkisi
ayrıntılı biçimde değerlendirilir.

\subsection{Harf Patlatma ve Yerçekimi Mekaniği}

Geçerli kelime silindikten sonra ilgili hücreler boşaltılır ve
üst hücrelerdeki harfler yerçekimi etkisiyle aşağıya kayar.
Boş kalan üst hücreler yeni harflerle doldurulur. Bu güncelleme,
gridin her hamleden sonra tutarlı ve oynanabilir kalmasını sağlar.

\subsection{Özel Güç Mekanikleri}

Özel güçler, kelime uzunluğuna bağlı olarak üretilir ve
kelimenin son harfinin bulunduğu hücrede saklanır. 4 harfli
kelimeler satır temizleme, 5 harfli kelimeler belirli bir
alanı patlatma, 6 harfli kelimeler sütun temizleme ve 7+
harfli kelimeler mega patlatma etkisi oluşturur. Bu güçler,
hücre tekrar seçildiğinde aktive edilir ve ilgili satır,
sütun veya alan üzerindeki harfleri temizler.

Özel güç aktarımı, hücreye bağlı bir durum olarak saklandığı
için, gücün kullanımı oyuncunun sonraki hamle planlamasına
bağlanır. Bu yaklaşım, kelime uzunluğunu teşvik ederken
tahta temizleme stratejilerinin bilinçli kurulmasını sağlar.

\subsection{Joker Sistemi}

Uygulamada, jokerlerin edinilmesi ve yönetilmesi için
ayrı bir market ekranı bulunur. Oyuncu, oyuna yüksek
miktarda altın bakiyesi ile başlar ve bu kaynak ile
farklı joker türlerini satın alabilir. Satın alınan
jokerler, Room tabanlı envanter yapısında veya yerel
veri modellerinde saklanır; böylece oyun yeniden
açıldığında miktarlar korunur.

Joker türleri; Balık (belirli bir harfi hedefleyerek
tekil temizleme), Tekerlek (rastgele bir bölgede
temizleme etkisi), Lolipop Kırıcı (seçilen satır
ve sütun kesitini patlatma), Serbest Değiştirme
(seçili bir harfi farklı bir harfe dönüştürme),
Harf Karıştırma (griddeki harflerin yeniden
dağıtılması) ve Parti Güçlendiricisi (tahta
genelinde geniş alan etkisi) olarak tanımlanmıştır.
Bu jokerler oyun ekranında seçildiğinde, hedefleme
modu etkinleşir; oyuncu ilgili hücre veya bölgeyi
işaretleyerek efekti tetikler. Hedefleme tamamlandığında
joker etkisi uygulanır, hamle akışı güncellenir ve
envanterdeki joker miktarı azaltılır.

\begin{figure}[htbp]
	\centering
	\includegraphics{market_ekrani.png}
	\caption{Market Ekranı}
	\label{fig:market}
\end{figure}

\subsection{Combo Mekaniği}

Combo mekaniği, ana kelime içinde oluşan anlamlı alt
kelimelerin ek puan getirmesi üzerinden kurgulanmıştır.
Seçilen harf dizisinde tespit edilen alt kelimeler, temel
kelime puanına eklenir; aynı alt kelime birden fazla kez
sayılmaz ve tekil olarak değerlendirilir. Bu yapı, oyuncuyu
daha uzun ve iç içe kelimeler kurmaya yönlendirir ve
puanlama stratejisini zenginleştirir.

\subsection{Skor Tablosu ve Veri Saklama}

Kullanıcı adı, DataStore üzerinde saklanarak uygulama
yeniden açıldığında kimlik bilgisinin kaybolması
engellenir. Oyun geçmişi ve skor kayıtları ise Room
veritabanında tutulur; böylece oyun oturumları arası
kalıcı veri sağlanır ve raporlama için tutarlı bir
kaynak elde edilir.

Skor tablosunda toplam oyun sayısı, en yüksek puan,
ortalama puan, toplam kelime sayısı, en uzun kelime ve
toplam süre gibi özet istatistikler sunulur. Ayrıca
geçmiş oyun kartları listelenerek her oturuma ait
puan, süre ve kelime verileri görüntülenir. Bu
yaklaşım, kullanıcı deneyimini kalıcı veri ile
ilişkilendirir ve oyuncunun ilerlemesini izlenebilir
hale getirir.

\begin{figure}[htbp]
	\centering
	\includegraphics{skor_tablosu_ekrani.png}
	\caption{Skor Tablosu Ekranı}
	\label{fig:skor_tablosu}
\end{figure}

\section{Algoritmalar ve Yalancı Kodlar}

Bu bölümde, Word Crush oyunundaki kritik akışlara ait yalancı kodlar
verilmektedir. Akışlar, oyunun başlatılması, kelime işleme
ve oyun içi destek mekaniklerinin uygulanmasını kapsar.

\subsection{Oyun Başlatma ve Grid Üretimi Akışı}

\noindent\textbf{Yalancı Kod:}
\begin{flushleft}
\ttfamily
BAŞLA \\
KullanıcıBilgisi $\leftarrow$ DataStoreOku \\
OyunGeçmişi $\leftarrow$ RoomOku \\
GridBoyutu $\leftarrow$ SeviyeSeçimineGöreBelirle \\
DÖNGÜ \\
\hspace*{1em} Grid $\leftarrow$ AğırlıklıHarfÜret(GridBoyutu) \\
\hspace*{1em} OynanabilirMi $\leftarrow$ KelimeVarMı(Grid) \\
\hspace*{1em} IF OynanabilirMi İSE ÇIK END IF \\
END DÖNGÜ \\
OyunDurumu $\leftarrow$ Başlat(Grid, HamleSayısı, Skor) \\
BİTİR
\end{flushleft}
\normalfont

\subsection{Kelime Seçimi ve Tahta Güncelleme Akışı}

\noindent\textbf{Yalancı Kod:}
\begin{flushleft}
\ttfamily
BAŞLA \\
Seçim $\leftarrow$ KomşubaşlıHarfSeç \\
IF SeçimUzunluğu $<$ 3 İSE \\
\hspace*{1em} HamleAzalt \\
\hspace*{1em} ÇIK
END IF \\
Kelime $\leftarrow$ SeçimiBirleştir \\
GeçerliMi $\leftarrow$ SözlükKontrol(Kelime) \\
HamleAzalt \\
IF GeçerliMi İSE \\
\hspace*{1em} PuanEkle(HarfPuanıToplam) \\
\hspace*{1em} ComboPuanEkle(AltKelimeler) \\
\hspace*{1em} SeçilenHücreleriTemizle \\
\hspace*{1em} HarfleriDüşür \\
\hspace*{1em} YeniHarfÜret \\
ELSE \\
\hspace*{1em} GeriBildirimGöster \\
END IF \\
BİTİR
\end{flushleft}
\normalfont

\subsection{Joker Kullanımı ve Skor Kaydı}

\noindent\textbf{Yalancı Kod:}
\begin{flushleft}
\ttfamily
BAŞLA \\
JokerSeç $\leftarrow$ MarkettenVeyaEnvanterdenSeç \\
Hedef $\leftarrow$ HücreVeyaBölgeİşaretle \\
JokerEtkisiUygula(Hedef) \\
GridGüncelle \\
OyunBittiMi $\leftarrow$ KontrolEt \\
IF OyunBitti İSE \\
\hspace*{1em} ÖzetİstatistikleriHesapla \\
\hspace*{1em} RoomSkorKaydet \\
END IF \\
BİTİR
\end{flushleft}
\normalfont

\section{Deneysel Sonuçlar}

Uygulama, Android emülatörü ve gerçek Android cihazlar üzerinde
çalıştırılmış; temel fonksiyonların ve oyun akışının beklendiği
şekilde çalıştığı doğrulanmıştır. Aşağıda kritik senaryolar
bağımsız alt başlıklar altında değerlendirilmiş, akışın doğruluğu
ve stabilitesi analiz edilmiştir.

\subsection{Kullanıcı Adı Kaydetme}

DataStore üzerinden kullanıcı adının kaydedildiği ve uygulama
yeniden başlatıldığında geri yüklenebildiği doğrulanmıştır.
Kimlik bilgisinin UI ve skor tablosu ekranlarında tutarlı
şekilde görüntülendiği gözlemlenmiştir.

\subsection{Yeni Oyun Başlatma ve Grid Oluşturma}

Yeni oyun başlatma aksiyonunun, belirlenen grid boyutu ile
tahta üretimini tetiklediği ve oynanabilirlik kontrolünden
geçen gridin ekrana yansıtıldığı test edilmiştir. Gridin
çözülmez durumlarda yeniden üretildiği ve kullanıcıya
geçerli bir tahta sunulduğu doğrulanmıştır.

\subsection{Kelime Doğrulama}

Seçilen harflerin sözlük veri yapısına göre doğrulandığı,
geçersiz kelimelerde hamlenin iptal edildiği ve geçerli
kelimelerde oyun akışının ilerlediği gözlemlenmiştir.
Minimum harf sayısı ve tekrarlı hücre kullanmama
kurallarının uygulandığı teyit edilmiştir.

\subsection{Puan Güncelleme}

Harf bazlı puanlama sisteminin doğru toplandığı, kelime
puanının skor tablosuna yansıtıldığı ve özel güç/ combo
çarpanlarının skoru beklenen biçimde etkilediği
görülmüştür.

\subsection{Combo Mekaniği}

Ana kelime içindeki anlamlı alt kelimelerin tespit edilip
ek puan ürettiği; aynı alt kelimenin tekrar sayılmadığı
senaryolarla doğrulanmıştır. Combo katkılarının skor
artışına tutarlı şekilde yansıdığı izlenmiştir.

\subsection{Özel Güçler}

4, 5, 6 ve 7+ harfli kelimelerle üretilen özel güçlerin
son harf konumunda saklandığı ve tekrar seçimde aktive
olduğu test edilmiştir. Satır/sütun temizleme ve
patlatma etkilerinin grid güncellemeleriyle uyumlu
çalıştığı gözlemlenmiştir.

\subsection{Joker Sistemi}

Market ekranından joker satın alma, envantere ekleme
ve oyun ekranında hedefleme ile kullanma akışları
test edilmiştir. Joker miktarının Room üzerinde
kalıcı olarak saklandığı ve kullanıldığında doğru
şekilde azaldığı görülmüştür.

\subsection{Skor Kaydı}

Oyun bitiminde özet istatistiklerin hesaplanarak Room
veritabanına kaydedildiği ve skor tablosu ekranında
geçmiş oyun kartlarıyla birlikte listelendiği doğrulanmıştır.
Bu akışın, kullanıcı deneyimini kalıcı veri ile
desteklediği değerlendirilmiştir.

\subsection{Teknik Zorluklar ve Değerlendirme}

Gelişme sürecinde, gesture yönetimi ve seçim izinin
doğru işlenmesi, Compose state yönetimi, grid güncelleme
senkronizasyonu ve veri kalıcılığı entegrasyonu gibi
teknik zorluklarla karşılaşılmıştır. Bu konular, UI ile
oyun motoru arasında net durum geçişleri tasarlanarak
çözülmüş; sonuç olarak uygulama beklenen oyun akışını
stabil biçimde sağlamıştır.

\section{Sonuç}

Bu çalışmada, Android tabanlı Word Crush mobil oyunu tasarlanmış ve
geliştirilmiş; oyun, belirlenen isterlerin büyük bir bölümünü karşılayan
stabil bir yapıya kavuşturulmuştur. Grid tabanlı oynanış, kelime
doğrulama, puanlama, combo, özel güç ve joker mekanikleri birlikte
çalışacak şekilde bütünleştirilmiş; kullanıcı verisi ve oyun geçmişi
kalıcı depolama katmanları ile desteklenmiştir.

Teknik açıdan; MVVM mimarisi, Jetpack Compose tabanlı arayüz, Room ve
DataStore tabanlı veri saklama, oyun motoru ve navigation akışları
başarıyla kurulmuş ve birbirleriyle uyumlu biçimde çalıştırılmıştır.
Bu yapının, mobil oyun geliştirme açısından modülerlik, test edilebilirlik
ve bakım kolaylığı sağladığı; kullanıcı deneyimi açısından ise
akıcı etkileşim, anlık geribildirim ve ilerleme izleme imkânlarını
artırdığı değerlendirilmiştir.

Gelecek çalışmalarda daha gelişmiş animasyonların eklenmesi, daha zengin
bir sözlük veri tabanının kullanılması, çevrim içi skor sistemi ile
rekabetçi özelliklerin artırılması ve farklı oyun modlarının
geliştirilmesi önerilmektedir. Bu iyileştirmelerin, oyunun hem
içerik zenginliğini hem de uzun vadeli oynanabilirliğini
artıracağı öngörülmektedir.

\vspace{1cm}

\section{Yazar Katkıları}

Bu proje, iki yazarın ortak çalışması ile geliştirilmiştir. Mobil
uygulama geliştirme süreci, oyun motoru tasarımı, veri saklama
katmanının kurulması, UI tasarımı, test faaliyetleri ve raporun
hazırlanması görevleri birlikte ve dengeli biçimde yürütülmüştür.
Çalışma boyunca kararlar ortaklaşmalı alınmış, uygulama ve dokümantasyon
aşamalarında eşit sorumluluk paylaşımı sağlanmıştır.

\vspace{1cm}

\begin{thebibliography}{12}

\bibitem{kotlin}
Kotlin, \emph{Kotlin Documentation}, [Çevrimici]. Erişim: \url{https://kotlinlang.org/docs/home.html}

\bibitem{compose}
Jetpack Compose, \emph{Compose Documentation}, [Çevrimici]. Erişim: \url{https://developer.android.com/jetpack/compose}

\bibitem{navigation}
Android Developers, \emph{Navigation Compose}, [Çevrimici]. Erişim: \url{https://developer.android.com/jetpack/compose/navigation}

\bibitem{room}
Room Persistence Library, \emph{Room Documentation}, [Çevrimici]. Erişim: \url{https://developer.android.com/jetpack/androidx/releases/room}

\bibitem{datastore}
Android Developers, \emph{DataStore}, [Çevrimici]. Erişim: \url{https://developer.android.com/topic/libraries/architecture/datastore}

\bibitem{material3}
Material Design 3, \emph{Material 3 in Compose}, [Çevrimici]. Erişim: \url{https://developer.android.com/jetpack/compose/designsystems/material3}

\bibitem{android}
Android Developers, \emph{Android Documentation}, [Çevrimici]. Erişim: \url{https://developer.android.com/}

\bibitem{gamedesign}
J. Schell, \emph{The Art of Game Design: A Book of Lenses}, 3rd ed. Boca Raton, FL, USA: CRC Press, 2019.

\bibitem{gamedesign2}
K. Salen and E. Zimmerman, \emph{Rules of Play: Game Design Fundamentals}. Cambridge, MA, USA: MIT Press, 2004.

\end{thebibliography}

\end{document}